\question[4]
Schreibe ein Python-Programm, das die Variablen \texttt{name = "Lena"} und
\texttt{punkte = 87} anlegt und diese mit \texttt{print(..., file=f)} in eine neue
Datei \texttt{ausgabe.txt} schreibt. Öffne die Datei dazu im Schreibmodus mit
\texttt{encoding='utf-8'} und schließe sie am Ende.

Erwarteter Inhalt der Datei \texttt{ausgabe.txt}:
\begin{lstlisting}
Lena hat 87 Punkte erreicht.
\end{lstlisting}

\begin{solutionbox}{4cm}
\begin{lstlisting}
name = "Lena"
punkte = 87

f = open('ausgabe.txt', encoding='utf-8', mode='w')
print(f'{name} hat {punkte} Punkte erreicht.', file=f)
f.close()
\end{lstlisting}
\end{solutionbox}